\documentclass[14pt,a4paper]{extarticle}

\usepackage[utf8]{inputenc}
\usepackage[english]{babel}
\usepackage{setspace}
\usepackage{geometry}

\geometry{
    left=25mm,
    right=15mm,
    top=20mm,
    bottom=20mm
}

\setstretch{1.5}

\begin{document}


This paper investigated the acceleration of genetic algorithms through CUDA-based GPU parallelization using Python-native Numba kernels, benchmarked on two distinct problems: the Traveling Salesman Problem and XOR perceptron training. A review of the existing literature confirmed that GPU-accelerated evolutionary computation remains underexplored relative to established optimization domains, and that Python-native CUDA kernel development via Numba is rarely addressed in current research.

\bigskip\textbf{Results.} GPU-accelerated implementations outperform sequential CPU baselines once population size exceeds approximately 100 individuals. Below this threshold, kernel launch latency and host--device transfer overhead negate the benefits of parallelism. Above it, speedups range from 1.15$\times$ to 14.56$\times$ depending on problem type and population size. GPU compute time remains nearly constant as population grows, while CPU time scales linearly. This is consistent with $O(1)$ vs.\ $O(n)$ per-generation complexity for parallel fitness evaluation. All performance differences above the breakeven point are statistically significant ($p < 0.05$, Cohen's $d > 0.8$). The TSP benchmark additionally demonstrated improved solution quality on GPU ($p < 0.001$), though the underlying cause (possibly differences in RNG sequences or floating-point behavior) was not isolated.

\bigskip\textbf{Hypothesis Evaluation.} The hypothesis is confirmed. CUDA parallelization delivers significant acceleration that scales with both population size and fitness function complexity, as predicted - speedups reach up to 14$\times$ on an entry-level mobile GPU. GPU-maintained populations also yield better solution quality. The dominant performance bottleneck is host--device data transfer rather than kernel computation, indicating that further speedup gains are achievable by keeping the full evolutionary loop GPU-resident.

\bigskip\textbf{Framework.} The experimental codebase was also refactored into a modular, publicly available Python library (Appendix~\ref{app:library}). Validation experiments confirmed statistical consistency with the original benchmarks and reproducibility. The library includes reusable components for fitness evaluation, selection, crossover, mutation, and GPU memory management that can be extended to new optimization problems.

\bigskip\textbf{Future Research.} Subsequent work could explore fully GPU-resident evolutionary loops (to eliminate per-generation host--device transfers), multi-objective optimization kernels, adaptive operator selection on GPU, higher-end hardware where the breakeven threshold is expected to shift lower. Support for variable-length genomes and hybrid CPU--GPU scheduling are natural next steps for the framework.

\end{document}
